\documentclass[12pt]{article}
\usepackage[utf8]{inputenc}
\usepackage[T1]{fontenc}
\usepackage[margin=1in]{geometry}
\usepackage{setspace}
\usepackage{xcolor}
\usepackage[colorlinks=true, urlcolor=blue!60!black, linkcolor=blue!60!black]{hyperref}
\setlength{\parindent}{0em}
\setlength{\parskip}{0.6em}
\renewcommand{\baselinestretch}{1.0}
\begin{document}
\begin{center}
{\LARGE \textit{Hans v. Louisiana}, 134 U.S. 1 (1890)\par}
\vspace{0.5em}
{\large Case Brief --- Two-Page Edition\par}
\end{center}

\section{Memory Jogger}
Although the Eleventh Amendment's text covers only suits by out-of-state or foreign citizens, state sovereign immunity bars a citizen from suing his own unconsenting State in federal court, even on a federal question (\href{https://supreme.justia.com/cases/federal/us/134/1/}{134 U.S. 1}).

\section{Detailed Case Facts}
Louisiana issued ``consolidated bonds'' in 1874 and pledged a tax to pay them, but its 1879 constitution remitted the coupons due January 1, 1880, and diverted the collected taxes. Bernard Hans, a Louisiana citizen holding those coupons, sued the State for \$87,500, claiming a Contract Clause violation and invoking federal-question jurisdiction.

\section{Procedural History}
The federal circuit court sustained Louisiana's objection that it could not be sued without consent and dismissed; on writ of error, the Supreme Court affirmed in 1890.

\section{Judicial Votes}
Affirmed, 9--0 (8--1 on reasoning). Bradley wrote for eight; Harlan concurred in the judgment only.

\section{Holding}
An unconsenting State cannot be sued in federal court by its own citizen, even under federal law, because the Constitution presupposes the States' preexisting sovereign immunity, which was not surrendered in the plan of the Convention.

\section{Analysis of Opinions}
Bradley rejected a literal reading of the Amendment as producing an ``anomalous result,'' since outsiders could not bring the same claims. \textit{Chisholm v. Georgia} had caused ``such a shock of surprise throughout the country'' that the Amendment was swiftly adopted, and Hamilton, Madison, and John Marshall had assured ratifiers that States could not be haled into federal court. The Judiciary Act of 1875 gave jurisdiction only ``concurrent with the courts of the several States,'' and Marshall's contrary dictum in \textit{Cohens v. Virginia} was ``extrajudicial.'' Harlan concurred only in the judgment, objecting that the attack on \textit{Chisholm} was unnecessary and that \textit{Chisholm} ``was based upon a sound interpretation of the Constitution as that instrument then was.''

\section{Examples: Future Applications}
\textbf{Same-Side:} a Texas resident's federal suit against Texas for overtime damages under a Commerce Clause statute is barred, as \textit{Seminole Tribe} confirmed. \textbf{Opposite-Side:} a suit against the Attorney General for a prospective injunction proceeds under \textit{Ex parte Young}. \textbf{Fence-Sitter:} Congress authorizing private suits against States under war or eminent-domain powers, which turns on immunity surrendered ``in the plan of the convention.''

\section{Critique}
Diversity theorists such as William Fletcher and Akhil Amar argue \textit{Hans} defies the text, which repealed only diversity jurisdiction; Baude, Sachs, and Nelson defend it as background common law. Progressive scholars see a post-Reconstruction retreat from Redeemer bond repudiation that leaves rights without remedies.

\section{Key Quotations}
\begin{quote}
``The suability of a State without its consent was a thing unknown to the law.''
\end{quote}
\noindent --- Justice Bradley, Opinion of the Court, 134 U.S. at 16.

\end{document}
